\begin{chapterbox}
    \chapter{Die Flamme von Choranat (WIP)}
    \label{wip-die-flamme-von-choranat}
    \az{Jahr 49}

    Die Runenmeisterin Burmrit, die Feuerkriegerin Falma, der Wolfskrieger Orfen und der Seher Leander stellen sich mysteriösen Spiegelwesen, die einem uralten Bannfeuer entsprangen.
\end{chapterbox}



Diese Fan-Geschichte ist noch in Arbeit. Ich teile sie bereits hier, damit sie für mich nicht verloren geht, falls mein Laptop spontan verschwinden sollte, und damit sie für euch nicht völlig verloren geht, falls ich spontan verschwinden sollte. Fühlt euch in letzterem Fall vollkommen frei, hierrin zu schmökern und nach Lust und Laune weiter daran zu werkeln. Aber solange ich ein aktiver Andori bleibe, lest bitte nicht zu genau nach, was ich geplant haben könnte. LG BBB :D





\newpage
\section{Prolog}
\az{Jahr 3}

Vier \gend{Magier*innen}{Magier} schufen den Bannkreis von Choranat. Ein Ritualkreis, der \gend{je*de}{jeden} von ihnen mit einem Spiegelbild konfrontierte, an dem sie wachsen konnten. Würden sie sich alle daran stärken können, würde das Feuer ihnen fortan eine Vervielfachung ihrer magischen Kräfte ermöglichen. Doch die \gend{Magier*innen}{Magier} wurden von ihren Spiegelbildern bezwungen. 

Der Bannkreis flackerte und verschwand.



\az{Jahr 26}

In einem anderen Land waren vier Personen zugegen. Da entzündete sich aus dem Nichts ein Bannfeuer. In den Flammen erkannten sie düstere Spiegelbilder ihrer selbst. Ängstlich rannten sie davon.

Das Bannfeuer loderte auf, bewegte sich durchs Land, und wuchs und wuchs und wuchs, bis es an einem Ort zur Ruhe kam. Ein Bannkreis von Choranat formte sich.









\newpage
\section{Kapitel}
\az{Jahr 49}

Der Kartograph Merrik und sein Beschützer, der Wolfskrieger Orfen, reisten durchs Steppenland der drei Brüder. Merrik hatte die Ausläufer des Grauen Gebirges rund um die Tote Festung kartographiert und wollte nun durch die Steppe zurück zum Wachsamen Wald gelangen.

Eine Zeit lang wurden die beiden begleitet von der Schamanin Asbark vom Iquar-Stamm und ihrem jugendlichen Lehrling Barz. Asbark teilte einen Sack voller Vorhersehungspulver mit Orfen, welcher ein wenig davon konsumierte und in der darauffolgenden Nacht Albträume von Flammen und einem dunklen Selbst hatte.


\trennzeichen


Die halb ertrunkene Feuerkriegerin Falma erreichte die Küste des Wachsamen Wald. Sie war eine blinde Passagierin auf einem Handelsschiff gewesen und zu früh über Bord gesprungen, um nicht entdeckt zu werden. 

Sie war eine Glutträgerin, kurz vor Beginn ihres letzten ihrer dreizehn Jahre, die sie das Glühende Herz des Flammenden Gottes in sich trug und außerhalb von Danwar nach dem Land der Tausend Flammen und dem flammenden Gott suchte. Seit dem Jahr 37 a.Z. war sie auf der Suche und 50 a.Z. würde das Glühende Herz ihren Körper zu sehr strapaziert haben und sie verbrennen.

Sie hatte eine Vision erhalten von einem Bannfeuer, welches nur von vier tapferen \gend{Held*innen}{Helden} aufgehalten werden konnte – jemand Feuriges, jemand mit weißen Augen, jemand einem leuchtenden Pulver und jemand mit einem Runenhammer.

Falma fand sich wieder in der Hütte des Sehers Leander, der ihre Ankunft vorhergesehen hatte und sie so gut wie möglich verarztete. Doch er konnte ihr Fieber nicht verringern.

Leander schleppte die verletzte Falma zur Hexe Reka und deren Schüler Meres. Erst die beiden konnten Falmas Fieber senken und sie aus dem Delirium reißen. 

Reka weigerte sich jedoch, Falma langfristig zu heilen. Die Bürde des Glühenden Herzens sei ihre Entscheidung gewesen. 

Meres übergab Leander einen knisternden Pulversack mit längerfristig heilendem Material, als Reka nicht hinsah.

Falma glaubte, in Leander die weißen Augen aus ihrer Vision zu erkennen. Leander erkannte in ihrer Erzählung eigene Visionen von Flammen.

Leander brachte Falma zum Baum der Lieder, um in den dortigen Archiven nachzuforschen.

In der Hütte des Obersten Bewahrers Melkart und dessen Nichte Josella waren gerade Merrik und Orfen zu Besuch. Merrik berichtete von seinen Abenteuern im Steppenland, was vor allem Melkarts und Josellas Adoptivkinder interessierte.

Falmas Erzählungen deckten sich mit Orfens Albträumen. Falma glaubte, in ihm die Person mit dem leuchtenden Pulver zu erkennen.

Falma, Orfen und Leander beschlossen, nach Silberhall aufzubrechen, wo man sich sowohl mit Runenmagie als auch mit dem Vertreiben von dunklen Geistern auskannte. Leander hielt nicht viel von Falmas und Orfens Interpretationen ihrer vagen Visionen, hätte aber durchaus Interesse daran, sein Runenwissen zu erweitern, war jenes in Cavern doch so gut gehütet. Und vor allem hatte er kürzlich zum ersten Mal seit Jahrzehnten Gerüchte gehört, Callem sei am Silberberg gesuchtet worden, auch wenn er das nicht mit den anderen teilte.

Merrik spielte mit dem Gedanken, sich ihnen anzuschließen, und eine Karte der Nebelinseln anzufertigen, begnügte sich aber am Ende damit, diesen Traum zu vertagen. Es gab so viel Material aus seiner letzten Reise, das er nachzubearbeiten hatte.

Auf der Reise bot Leander Falma das heilende Pulver an. Es sollte ihr eine Chance geben, die 13 Jahre zu überleben. Falma lehnte fanatisch ab.

In Silberland schloss sich die Runenmeisterin Burmrit der Gruppe an. Falma glaubte, in ihr die Person mit dem Runenhammer zu erkennen.

Gemeinsam spürte das Quartett am dem Silberberg das Bannfeuer an, welches umherloderte und wuchs und wuchs. Die vier stellten sich dem Feuer, statt davonzurennen. Es formte sich ein Bannkreis von Choranat und die vier wurden mit ihren Spiegelbildern konfrontiert.

Burmrit versuchte, die Spiegelbilder von Choranat wie Geister in Edelsteine zu bannen, war allerdings nicht erfolgreich. Leander steckte einen dieser Edelsteine ein – es könnte in Zukunft äußerst nützlich sein, zu wissen, mit welchen Runen man Geister in Stein bannte.

Ein Kampf gegen die verschiedenen Spiegelbilder entbrannte. 

Leander und sein Spiegelbild verloren sich zu einem geistigen Kampf der Gewitztheit und Vorhersehung.

Burmrit konnte ihr Spiegelbild im letzten Moment von Varatans Fluch erwischen und in den Norden ziehen lassen.

Orfen und Falma besiegten im direkten Kampf ihre Spiegelbilder, doch nicht ehe Falma ihr Leben opferte, um Orfens zu retten.

Nach Falmas Tod glühte das Glühende Herz auf und verkohlte ihren Leichnam zu Asche. Das Glühende Herz blieb zurück darin.


\newpage
\section{Epiloge}

Orfen und Leander besuchten Danwar. Der eine, um das Glühende Herz und Falmas Asche zu ihrer Familie zu bringen, der andere, um diese legendäre Insel zum ersten Mal selbst zu erleben. Das Glühende Herz konnte als Kompass dienen.

Sie wurden nicht freundlich empfangen. Als versagt habende Glutträgerin wurde Falmas Name aus den Archiven gestrichen und nicht gerne über sie gesprochen. Und das Glühende Herz wäre offenbar von selbst wieder zurück auf Danwar erschienen, wenn die 13 Jahre ausgelaufen waren.

Orfen ließ das Vorhersehungspulver von Asbark und das heilende Pulver vor Meres zurück. Der grantige Kord nahm viel zu viel dieser Mischung ein. Leander nahm mit Horror war, wie Kords Sehkraft ungeahnte Macht einnahm und Kord dadurch litt.


\trennzeichen

Burmrit kehrte zurück und unterhielt sich mit ihrer Schülerin Iril. Sie erzählte ihr von der Prophezeiung, dass der Bannkreis von Charant aufgehalten würde jemand Feurigem, jemand mit weißen Augen, jemand mit leuchtenden Pulvern und jemand mit einem Runenhammer.



\trennzeichen

Orfen beendete seine Erzählung dieser Ereignisse gegenüber Merrik, welche sie zu Pergament brachte und in den Untiefen der Archive des Baum des Lieder verstaute.


\az{Jahr 72}

Vier wohlbekannte \gend{Bewahrer*innen}{Bewahrer} waren im Wachsamen Wald am Beeren sammeln, als sich zwischen ihnen aus dem Nichts ein Bannfeuer entzündete. Sie sahen düstere Spiegelbilder von sich auf der anderen Seite des Feuers und zogen sich zurück. Auch das Feuer raste davon, und die Flamme wuchs und wuchs.

\trennzeichen

Der Bewahrer Phlegon, still ängstlich von der Begegnung mit dem Bannfeuer, brachte Merriks Aufzeichnungen wieder ans Licht. Im Kerzenschein setzte er sich an einen Tisch, schlürfte seinen heißen Cocoa und studierte aufmerksam Merriks Zeichnungen vom Bannkreis von Choranat.




